\documentclass[10pt,a4paper]{amsart}
\usepackage{amsmath,amssymb,amsthm,graphicx}
\usepackage{verbatim}
\newenvironment{reference}{\comment}{\endcomment}
%\newenvironment{reference}{}{}
\newenvironment{slogan}{\comment}{\endcomment}
\newenvironment{history}{\comment}{\endcomment}

% For commutative diagrams we use Xy-pic
\usepackage[all]{xy}

% We use 2cell for 2-commutative diagrams.
\xyoption{2cell}
% Theorem environments.
%
\theoremstyle{plain}
\newtheorem{theorem}[subsection]{Theorem}
\newtheorem{proposition}[subsection]{Proposition}
\newtheorem{lemma}[subsection]{Lemma}

\theoremstyle{definition}
\newtheorem{definition}[subsection]{Definition}
\newtheorem{example}[subsection]{Example}
\newtheorem{exercise}[subsection]{Exercise}
\newtheorem{situation}[subsection]{Situation}

\theoremstyle{remark}
\newtheorem{remark}[subsection]{Remark}
\newtheorem{remarks}[subsection]{Remarks}

\numberwithin{equation}{subsection}

% Macros
%
\def\lim{\mathop{\mathrm{lim}}\nolimits}
\def\colim{\mathop{\mathrm{colim}}\nolimits}
\def\Spec{\mathop{\mathrm{Spec}}}
\def\Hom{\mathop{\mathrm{Hom}}\nolimits}
\def\Ext{\mathop{\mathrm{Ext}}\nolimits}
\def\SheafHom{\mathop{\mathcal{H}\!\mathit{om}}\nolimits}
\def\SheafExt{\mathop{\mathcal{E}\!\mathit{xt}}\nolimits}
\def\Sch{\mathit{Sch}}
\def\Mor{\mathop{\mathrm{Mor}}\nolimits}
\def\Ob{\mathop{\mathrm{Ob}}\nolimits}
\def\Sh{\mathop{\mathit{Sh}}\nolimits}
\def\NL{\mathop{N\!L}\nolimits}
\def\CH{\mathop{\mathrm{CH}}\nolimits}
\def\proetale{{pro\text{-}\acute{e}tale}}
\def\etale{{\acute{e}tale}}
\def\QCoh{\mathit{QCoh}}
\def\Ker{\mathop{\mathrm{Ker}}}
\def\Im{\mathop{\mathrm{Im}}}
\def\Coker{\mathop{\mathrm{Coker}}}
\def\Coim{\mathop{\mathrm{Coim}}}

% Boxtimes
%
\DeclareMathSymbol{\boxtimes}{\mathbin}{AMSa}{"02}

%
% Macros for moduli stacks/spaces
%
\def\QCohstack{\mathcal{QC}\!\mathit{oh}}
\def\Cohstack{\mathcal{C}\!\mathit{oh}}
\def\Spacesstack{\mathcal{S}\!\mathit{paces}}
\def\Quotfunctor{\mathrm{Quot}}
\def\Hilbfunctor{\mathrm{Hilb}}
\def\Curvesstack{\mathcal{C}\!\mathit{urves}}
\def\Polarizedstack{\mathcal{P}\!\mathit{olarized}}
\def\Complexesstack{\mathcal{C}\!\mathit{omplexes}}
% \Pic is the operator that assigns to X its picard group, usage \Pic(X)
% \Picardstack_{X/B} denotes the Picard stack of X over B
% \Picardfunctor_{X/B} denotes the Picard functor of X over B
\def\Pic{\mathop{\mathrm{Pic}}\nolimits}
\def\Picardstack{\mathcal{P}\!\mathit{ic}}
\def\Picardfunctor{\mathrm{Pic}}
\def\Deformationcategory{\mathcal{D}\!\mathit{ef}}
\usepackage{cleveref}
\crefname{originalreference}{original reference}{original references}
\Crefname{originalreference}{Original reference}{Original references}
\makeatletter
\newlabel{definition-rig-smooth-homomorphism}{{0GAI}{source}{}{sourcepointer-1}{}}
\newlabel{definition-rig-smooth-homomorphism@cref}{{[originalreference][0][]0GAI}{[1][0][]source}{}{}{}}
\makeatother
\crefname{originalsection}{original section}{original sections}
\Crefname{originalsection}{Original section}{Original sections}
\makeatletter
\newlabel{formal-spaces-section-morphisms-rings}{{0ANA}{source}{}{sourcepointer-2}{}}
\newlabel{formal-spaces-section-morphisms-rings@cref}{{[originalsection][0][]0ANA}{[1][0][]source}{}{}{}}
\makeatother
\crefname{originalreference}{original reference}{original references}
\Crefname{originalreference}{Original reference}{Original references}
\makeatletter
\newlabel{lemma-finite-type}{{04U5}{source}{}{sourcepointer-3}{}}
\newlabel{lemma-finite-type@cref}{{[originalreference][0][]04U5}{[1][0][]source}{}{}{}}
\makeatother
\crefname{originalreference}{original reference}{original references}
\Crefname{originalreference}{Original reference}{Original references}
\makeatletter
\newlabel{remark-NL-well-defined-topological}{{0GBX}{source}{}{sourcepointer-4}{}}
\newlabel{remark-NL-well-defined-topological@cref}{{[originalreference][0][]0GBX}{[1][0][]source}{}{}{}}
\makeatother
\title{Rig-smoothness over one ideal of definition implies all}
\author{The Stacks Project Authors}
\date{stacks-project:restricted.tex:local-snapshot-20261012; GNU FDL 1.2 or later}
\begin{document}
\maketitle
\noindent\textit{Editorial note.} All delivered mathematical text is the author's own, copied from the pinned Stacks Project chapter file restricted.tex; no statement or proof step was rewritten, reordered or removed.\par\medskip
\noindent\textit{Editorial note.} The counted proof is the author's complete argument: the counted statement is the author lemma environment labelled lemma-rig-smooth-continuous; the proof environment immediately following it is its own proof\par\medskip
\noindent\textit{Editorial note.} The label lemma-rig-smooth-continuous names the counted statement, and the proof environment immediately following it in the source is delivered verbatim as its proof.\par\medskip
\noindent\textit{Editorial note.} The counted statement is carried by the author's own lemma environment in the chapter source.\par\medskip
\noindent\textit{Allowed standard prerequisites (editorial).}\begin{itemize}
\item ideals of definition of a Noetherian adic ring are pairwise commensurable: for any two ideals of definition I and I' there is c >= 0 with I\textasciicircum{}c subset I' and (I')\textasciicircum{}c subset I
\item the definition of rig-smooth homomorphisms of WAdm\textasciicircum{}Noeth via the completed naive cotangent complex (definition-rig-smooth-homomorphism)
\item the completed naive cotangent complex NL\textasciicircum{}wedge\_\{B/A\} is independent of the chosen ideal of definition of A (remark-NL-well-defined-topological of the same chapter)
\end{itemize}\medskip
\ifdefined\hypertarget\hypertarget{sourcepointer-1}{}\fi
\noindent\textit{Source reference.} Original reference 0GAI: sources/stacks-restricted/pinned-original-source.tex, source lines 831--831. Prerequisite pointer invoked at proof line(s) 4233; it names a standard fact of the same chapter that the counted argument uses, and every step of the counted proof built on it is retained verbatim in the delivered body.\par\medskip
\ifdefined\hypertarget\hypertarget{sourcepointer-2}{}\fi
\noindent\textit{Source reference.} Original section 0ANA: sources/stacks-chow/tags.txt, source lines 13433--13433. Prerequisite pointer invoked at proof line(s) 4216; it names a standard fact of a sibling chapter that the counted argument uses, and every step of the counted proof built on it is retained verbatim in the delivered body.\par\medskip
\ifdefined\hypertarget\hypertarget{sourcepointer-3}{}\fi
\noindent\textit{Source reference.} Original reference 04U5: sources/stacks-restricted/pinned-original-source.tex, source lines 2439--2439. Prerequisite pointer invoked at proof line(s) 4220; it names a standard fact of the same chapter that the counted argument uses, and every step of the counted proof built on it is retained verbatim in the delivered body.\par\medskip
\ifdefined\hypertarget\hypertarget{sourcepointer-4}{}\fi
\noindent\textit{Source reference.} Original reference 0GBX: sources/stacks-restricted/pinned-original-source.tex, source lines 2487--2487. Prerequisite pointer invoked at proof line(s) 4237; it names a standard fact of the same chapter that the counted argument uses, and every step of the counted proof built on it is retained verbatim in the delivered body.\par\medskip
\begin{lemma}
\label{lemma-rig-smooth-continuous}
Let $A \to B$ be a morphism in $\textit{WAdm}^{Noeth}$
(Formal Spaces, Section \ref{formal-spaces-section-morphisms-rings}).
The following are equivalent:
\begin{enumerate}
\item[(a)] $A \to B$ satisfies the equivalent conditions of
Lemma \ref{lemma-finite-type} and there exists an ideal of definition
$I \subset B$ such that $B$ is rig-smooth over $(A, I)$, and
\item[(b)] $A \to B$ satisfies the equivalent conditions of
Lemma \ref{lemma-finite-type} and for all ideals of definition
$I \subset A$ the algebra $B$ is rig-smooth over $(A, I)$.
\end{enumerate}
\end{lemma}

\begin{proof}
Let $I$ and $I'$ be ideals of definitions of $A$. Then there exists an
integer $c \geq 0$ such that $I^c \subset I'$ and $(I')^c \subset I$. Hence
$B$ is rig-smooth over $(A, I)$ if and only if
$B$ is rig-smooth over $(A, I')$. This follows from
Definition \ref{definition-rig-smooth-homomorphism},
the inclusions $I^c \subset I'$ and $(I')^c \subset I$, and
the fact that the naive cotangent complex $\NL_{B/A}^\wedge$
is independent of the choice of ideal of definition of $A$ by
Remark \ref{remark-NL-well-defined-topological}.
\end{proof}

\end{document}
